\section{Related Work}
\label{sec:related-work}

Random finite-set filters model uncertain cardinality and association. PHD
filtering gives a first-moment recursion \cite{Mahler2003PHD}; labeled RFS/LMB
models retain identity \cite{Vo2014LRFS,Reuter2014LMB}, and multi-sensor variants
support cooperation \cite{Vo2019MSGLMB}. Under unknown cross-correlations,
KLA/GCI forms a normalized geometric mean \cite{Battistelli2014KLA}, with PHD
and multi-Bernoulli RFS formulations \cite{Uney2013EMDPHD,Wang2017MBGCI}.
Consensus and robust LMB/LRFS methods address networked tracking, including
label-set inconsistency
\cite{Fantacci2015ConsensusLRFS,Shen2022ConsensusLMB,
Li2018RobustDistributedLRFS,Li2019EfficientMultiAgentLRFS}. Minimum-information-
loss and arithmetic-average LMB rules instead use different density-level
objectives \cite{Gao2020MIL,Wei2024AALMB}. These works determine how labeled
posteriors are fused. We propose no new density-level rule and audit only the
custom presence-conditioned per-label update defined below.

\newpage

Communication methods alter a different layer. Event-triggered LMB methods
change when updates are transmitted
\cite{Shen2022EventTriggeredLMB,Li2023EventTriggeredConsensusLMB}; related
triggers act on $M\delta$-GLMB labels or hypotheses
\cite{Li2026EventTriggeredMdeltaGLMB}. Partial-consensus PHD changes which GM
components are disseminated \cite{Li2019PartialConsensusGMPHD}. State and
covariance projection constructs compact track-level inputs for asynchronous
Gaussian fusion \cite{Xue2026StateCovProjection}. In contrast, we hold the
fusion rule, schedule, selection, and receiver fixed, then derive the message
from that receiver's first projection. The certificate covers label/existence,
Gaussian overlap, missing-label normalization, and the codec. Triggering could
be composed with this interface, but that combination is not evaluated here.
